\documentclass{article}

\usepackage[a4paper, margin=1in]{geometry}
\usepackage[T1]{fontenc}
\usepackage[utf8]{inputenc}
\usepackage{times}
\usepackage{amsmath,amssymb,amsthm}
\usepackage{booktabs}
\usepackage{array}
\usepackage{setspace}
\usepackage{microtype}
\usepackage[numbers,compress]{natbib}
\usepackage{xcolor}
\usepackage{titlesec}
\usepackage{parskip}
\usepackage{enumitem}
\usepackage[colorlinks=true,linkcolor=black,citecolor=black,urlcolor=blue]{hyperref}

\titleformat{\section}{\large\bfseries}{\thesection.}{0.5em}{}
\titleformat{\subsection}{\normalsize\bfseries}{\thesubsection}{0.5em}{}
\titleformat{\subsubsection}{\normalsize\bfseries\itshape}{\thesubsubsection}{0.5em}{}

\setlength{\parskip}{6pt}
\setlength{\parindent}{0pt}
\onehalfspacing

\newtheorem{theorem}{Theorem}
\newtheorem{definition}{Definition}
\newtheorem{proposition}{Proposition}

\title{\textbf{Acoustic Formant Features Are Insufficient for Semantic\\
Axis Recovery but Sufficient for Articulatory Locus Classification\\
in Sanskrit Verbal Roots: A Spiking Reservoir Computing Investigation}}

\author{Amit Kumar\\
\small Independent Researcher, Bihar, India\\
\small \href{mailto:c7ul6g@gmail.com}{c7ul6g@gmail.com}\\[4pt]
\small ORCID: \href{https://orcid.org/0009-0003-3362-5461}{0009-0003-3362-5461}}

\date{%
  April 2026\\[4pt]
  \small\textit{Preprint — Not yet peer reviewed}\\[2pt]
  \small\textcopyright\ 2026 Amit Kumar. Licensed under
  \href{https://creativecommons.org/licenses/by/4.0/}{CC BY 4.0}%
}

\begin{document}
\maketitle

\begin{center}
\small\textit{This paper reports a planned negative result arising from
a systematic audit of the DDIN research program (Papers 1--3). The audit
identified a labeling circularity in earlier work; this paper presents
the corrected empirical findings and their implications.}
\end{center}

%% ═══════════════════════════════════════════════════════════════════
\begin{abstract}
The phonosemantic hypothesis holds that in motivated sign systems such
as Sanskrit, the acoustic phonological form of a verbal root should carry
statistically detectable information about its semantic category. We test
this hypothesis using an Adaptive Exponential Leaky Integrate-and-Fire
(AdEx) spiking reservoir with per-neuron BCM homeostasis, processing the
full 2,000-root Pāṇinian \textit{Dhātupāṭha} under three controlled
input conditions: acoustic features only, semantic prior only, and full
acoustic-semantic fusion.

Under a rigorous labeling regime derived from Monier-Williams dictionary
glosses with word-boundary substring matching (coverage: 76.2\%,
1,722~roots), we find: (A)~acoustic-only ARI $= 0.013$, indistinguishable
from chance; (B)~prior-only ARI $= 0.70$; (C)~full-fusion ARI $= 0.69$,
representing a negative marginal acoustic contribution ($\Delta = -0.01$).
A Speed-1 experiment confirms ARI $= 0.016$ at the initial consonant.

A Phase~13 architectural validation experiment rules out reservoir failure
as the cause: the reservoir recovers the Pāṇinian articulatory locus
taxonomy (five consonant place-of-articulation classes) at ARI $= 0.0494$
($p < 0.001$). A semantic-locus correlation experiment finds ARI $= 0.0095$
between the Monier-Williams semantic taxonomy and the locus taxonomy,
quantifying their near-orthogonality directly. The reservoir fails to
recover the semantic taxonomy not because it cannot process acoustic
information, but because the target taxonomy does not align with the
articulatory structure the acoustic features encode.

These findings identify the failure mechanism precisely: first-order
acoustic formants are insufficient for the Monier-Williams semantic
taxonomy, but sufficient for articulatory locus classification. The
architecture is functional; the task-feature mismatch is the binding
constraint. Two engineering contributions stand independently: the
Epileptiform Synchrony Limit characterization~\citep{kumar2026esl} and
the BCM homeostasis architecture that resolves it.

\smallskip
\noindent\textbf{Keywords:} phonosemantics, Sanskrit, spiking neural
networks, reservoir computing, AdEx, BCM homeostasis, negative result,
Dhātupāṭha, acoustic insufficiency, articulatory locus, motivated sign
\end{abstract}

%% ═══════════════════════════════════════════════════════════════════
\section{Introduction}

\subsection{The Phonosemantic Hypothesis}

The hypothesis of phonosemantic correspondence holds that in certain
languages, the acoustic form of a word is non-arbitrarily related to its
meaning \citep{hinton1994,magnus2001,blasi2016}. This position challenges
the Saussurean principle of the arbitrariness of the linguistic sign
\citep{saussure1916}, which underlies virtually all modern statistical
language modeling. Sanskrit has been proposed as a language in which
this correspondence has been most rigorously formalized through the
Pāṇinian grammatical tradition \citep{cardona1997}: the
\textit{Dhātupāṭha}, a catalog of approximately 2,000 verbal roots
classified by semantic class (\textit{gaṇa}), provides the experimental
substrate against which this claim can be tested computationally.

The DDIN (Devavāṇī-Derived Interpretable Network) research program
\citep{kumar2026phonosemantic,kumar2026ddin,kumar2026esl} was designed
to test whether acoustic phonological features derived from locus theory
and the International Phonetic Alphabet can organize Sanskrit verbal roots
into their traditional semantic classes without semantic supervision. Papers
1 and 2 of the program reported above-chance Adjusted Rand Index (ARI)
results on a 150-root benchmark, and Paper 4 reported ARI $= 0.9758$ on
the full 2,000-root corpus.

\subsection{The Labeling Audit}

A systematic audit of the experimental code underlying Papers 2 and 4
identified a structural circularity in the semantic axis label assignment
function \texttt{extract\_artha\_stem\_tensor()}. In the versions used
to generate those results, labels were assigned based on the initial
consonant of the verbal root (\texttt{root\_slp1[0].lower()}), mapping
velars to Motion, palatals to Experiential, retroflex to Transformation,
dentals to Separation, and labials to Containment. The acoustic feature
vector was simultaneously derived from the same initial consonant via
locus equations and formant tables. The ARI therefore measured the
network's ability to separate \emph{consonant place classes} from each
other---a task trivially solved by the input encoding itself---rather than
genuine acoustic-to-semantic recovery.

This paper reports the results of re-running the key experiments with a
corrected labeling function that derives semantic axis membership
exclusively from dictionary gloss keywords matched against the
Monier-Williams Sanskrit-English Dictionary \citep{monierWilliams1899}
using word-boundary tokenization. The corrected function is described in
Section~\ref{sec:methods}.

\subsection{Scope of This Paper}

This paper makes four claims:

\begin{enumerate}[noitemsep]
  \item Under honest gloss-based labeling, acoustic-only ARI is 0.013
    (chance level) across multiple acoustic encoding schemes.
  \item Full acoustic-semantic fusion ARI (0.69) is \emph{lower} than
    prior-only ARI (0.70), indicating the acoustic channel contributes
    noise, not signal.
  \item The Speed-1 finding of Paper 2 (ARI $= 0.20$ at first phoneme)
    does not replicate under corrected labels (ARI $= 0.016$).
  \item The BCM homeostasis engineering contribution and the ESL
    characterization \citep{kumar2026esl} are independent of the
    labeling regime and remain valid.
\end{enumerate}

%% ═══════════════════════════════════════════════════════════════════
\section{Methods}
\label{sec:methods}

\subsection{Corpus and Labeling}

The corpus is the full \textit{Dhātupāṭha} sourced from the
\texttt{ashtadhyayi.com} dataset \citep{goyal2015}, providing 2,259
verbal roots with Devanagari transliteration and gloss in Sanskrit.
Roots are transliterated to SLP1 encoding for processing.

\textbf{Corrected label assignment.} For each root, the semantic axis
label is derived from the SLP1 transliteration of the dictionary gloss
(\textit{artha}) using word-boundary tokenization:

\begin{verbatim}
tokens = re.split(r'[\s,;।]+', artha_slp1.lower())
for stem in axis_stems[axis]:
    if any(t.startswith(stem) for t in tokens):
        tensor[axis] = 1.0
\end{verbatim}

Stems are drawn from a 300-entry Monier-Williams derived vocabulary
covering the five semantic axes (Motion: \textit{gat, cal, gam, car,
kram, \ldots}; Separation: \textit{hiMs, Bed, Cid, nAS, \ldots};
Containment: \textit{Dar, pAl, rakz, banD, \ldots}; Experiential:
\textit{dfS, Sru, jYA, man, \ldots}; Transfer: \textit{dA, veza, yAc,
niD, \ldots}). Roots with no matched stem receive the uniform prior
$[0.2, 0.2, 0.2, 0.2, 0.2]$ and are excluded from labeled evaluation.
Coverage under this scheme is 76.2\% (1,722 of 2,259 roots).

\textbf{Ablation baseline.} K-Means ($k=5$, 100 restarts) applied
directly to the 5D prior vectors achieves ARI $= 0.70$, establishing
the gloss-label floor. An MLP classifier (two hidden layers, 64 and 32
units, 5-fold cross-validation) achieves ARI $= 1.0$ on the prior
vectors, confirming perfect linear separability of the gloss-based
taxonomy.

\subsection{Acoustic Feature Encoding (23D)}

Each root is encoded as a 23-dimensional acoustic vector by
concatenation of three feature groups:

\textbf{Locus equations (6D).} $F_1$ and $F_2$ formant transition
coefficients for the initial consonant, derived from IPA locus equation
tables \citep{sussman1991}. Individual phoneme-specific values are used
within place classes (e.g., /p/ at $(F_1, F_2) = (0.35, 0.75)$,
/b/ at $(0.38, 0.72)$, /m/ at $(0.42, 0.68)$) to provide
within-class variation absent from the original class-level encoding.

\textbf{Phonation profile (5D).} Binary features encoding voicing,
aspiration, nasality, sibilance, and approximant status of the initial
consonant.

\textbf{Steady-state formants (12D).} $F_1$ and $F_2$ values for up
to six phoneme positions, zero-padded for shorter roots.

\subsection{Reservoir Architecture}

The reservoir is a two-layer AdEx spiking neural network (MegaSNN):
Layer 1 consists of 512 AdEx neurons ($a = 0$, $b = 0$, $\tau_w = 5$\,ms,
$g_L = 5$\,S/m$^2$) with a sparse input projection (18\% connectivity,
weight scale 2500). Layer 2 consists of 256 AdEx neurons ($a = 2$,
$b = 10$, $\tau_w = 30$\,ms, $g_L = 2$\,S/m$^2$). Per-neuron BCM
homeostasis operates on Layer 2 with $\tau_\theta = 500$\,ms. No
recurrent weights are trained; $W = 0$ throughout (Receiver Model
\citep{kumar2026ddin}).

Each root is presented for 20 timesteps followed by 20 timesteps of
silence. Reservoir states are recorded as Layer 2 membrane potentials
$V_i$ at stimulus offset, $\ell_2$-normalized prior to clustering.

\textbf{Clustering and ARI.} K-Means ($k = 5$, 100 restarts,
$\ell_2$-normalized 256D state space) is applied to the recorded states.
ARI is computed between cluster assignments and gloss-based labels using
\texttt{sklearn.metrics.adjusted\_rand\_score}. $k = 5$ is fixed
\textit{a priori} and not tuned against ARI.

\subsection{Three-Condition Ablation}

\begin{description}[noitemsep]
  \item[Condition A (Acoustic-only).] Reservoir input:
    $[\mathbf{a}_r \,\|\, \mathbf{0.2}_{5}]$. Acoustic features with
    uniform prior encoding maximum semantic uncertainty.
  \item[Condition B (Prior-only).] Reservoir input:
    $[\mathbf{0}_{23} \,\|\, \boldsymbol{\alpha}_r]$. Zero acoustic
    channel; gloss-based prior vector only.
  \item[Condition C (Full fusion).] Reservoir input:
    $[\mathbf{a}_r \,\|\, \boldsymbol{\alpha}_r]$. Both channels active.
\end{description}

\subsection{Speed-1 Experiment}

To test whether the initial consonant carries detectable semantic signal
under corrected labels, we run Condition A with truncated input: roots
are presented as single-phoneme sequences (initial consonant only, then
silence). ARI is measured against gloss-based labels.

%% ═══════════════════════════════════════════════════════════════════
\section{Results}

\subsection{Three-Condition Ablation}

\begin{table}[h]
\centering
\begin{tabular}{llcl}
\toprule
Condition & Input & ARI & Interpretation \\
\midrule
MLP baseline & 5D prior only & 1.000 & Perfect separability of gloss labels \\
B: Prior-only & $[\mathbf{0}_{23} \| \boldsymbol{\alpha}]$ & 0.702 &
  Reservoir degrades prior signal \\
C: Full fusion & $[\mathbf{a} \| \boldsymbol{\alpha}]$ & 0.689 &
  Acoustic channel adds noise ($\Delta = -0.013$) \\
A: Acoustic-only & $[\mathbf{a} \| \mathbf{0.2}]$ & 0.013 &
  Chance level \\
\bottomrule
\end{tabular}
\caption{Three-condition ablation results under corrected gloss-based
  labeling ($N = 1{,}722$ matched roots, $k = 5$, 100 K-Means restarts).
  Mean firing rate: 2.8 spikes/root (stable, no ESL observed).}
\label{tab:ablation}
\end{table}

The acoustic-only condition achieves ARI $= 0.013$, indistinguishable
from chance ($p > 0.05$ under permutation test, $n = 1000$ permutations).
This result is stable across two acoustic encoding schemes: the original
consonant-class-level formants and the refined phoneme-specific formants
with within-class variation. Adding within-class variation increased
acoustic-only ARI from 0.010 to 0.013, a difference not significant and
insufficient to support the phonosemantic hypothesis.

The negative fusion delta ($\Delta = -0.013$) indicates the acoustic
features disrupt rather than augment the gloss-based semantic organization.
The reservoir projects the 28D input into a 256D spike-count manifold;
under random sparse projection, the 23D acoustic component adds
high-dimensional noise that reduces the separability of the 5D prior
signal.

The MLP baseline (ARI $= 1.0$) confirms the gloss-based labels are
perfectly linearly separable, ruling out label noise as an explanation
for the acoustic channel's failure.

\subsection{Speed-1 Experiment}

Under corrected gloss-based labels, ARI at the initial consonant is
$0.016$ (chance level). This does not replicate the Speed-1 finding
reported in Paper 2 (ARI $= 0.204$ at first phoneme). The earlier
result was an artifact of the consonant-based labeling regime: because
labels were assigned from the initial consonant and acoustic features
were derived from the same consonant, single-phoneme presentation was
sufficient to recover the label with high fidelity.

Under honest labels, no detectable semantic signal is present in the
initial consonant alone, nor in any longer integration window tested
(2, 3, 5, and 10 phonemes all yielded ARI $< 0.02$).

\subsection{Reservoir Behavior}

The BCM homeostasis architecture successfully maintained stable
asynchronous firing at 2.8 spikes/root throughout all experiments.
No epileptiform synchrony was observed. This confirms the engineering
contribution of the ESL resolution \citep{kumar2026esl} independent
of the semantic clustering failure: the reservoir is biophysically
stable; it simply does not carry the semantic signal of interest.

\subsection{Label Coverage and Quality}

Of 2,259 corpus roots, 1,722 (76.2\%) received a matched gloss-based
label. The 22.8\% unmatched roots received the uniform fallback prior
and were excluded from ARI evaluation. Manual inspection of 50 randomly
sampled matched roots confirmed correct axis assignment in 44 cases
(88\% precision). The six errors were cases of polysemous roots whose
primary gloss matched multiple axes; in all cases the dominant axis was
correctly assigned.

%% ═══════════════════════════════════════════════════════════════════
\section{Phase 13: Architectural Validation and Taxonomy Diagnosis}
\label{sec:phase13}

The three-condition ablation in Section~3 established that the reservoir
fails to recover the Monier-Williams semantic taxonomy from acoustic features.
This left open a critical ambiguity: is the failure located in the
\emph{architecture} (the reservoir cannot process acoustic sequences) or
in the \emph{task} (the semantic taxonomy is not recoverable from acoustic
features regardless of architecture)? Phase~13 was designed to resolve
this ambiguity using two targeted experiments.

\subsection{Experiment 56: Locus Taxonomy Baseline}

\textbf{Design.} We replaced the Monier-Williams semantic taxonomy with
the Pāṇinian articulatory locus taxonomy: the five place-of-articulation
classes that Sanskrit grammar assigns to every consonant (kaṇṭhya/velar,
tālavya/palatal, mūrdhanya/retroflex, dantya/dental, oṣṭhya/labial). This
taxonomy is defined by the initial consonant of each root and is therefore
directly encoded in the F2 locus equations of the 23D acoustic feature
vector. If the reservoir functions correctly, it should recover this
taxonomy above chance because the relevant signal is present in the input
by construction.

Labels were assigned deterministically from the SLP1 transliteration of
each root's initial consonant using an explicit consonant-to-class mapping.
\textbf{Note}: unlike Phase 12's label leakage bug, this labeling is openly
derived from the initial consonant and is stated as such. The locus
experiment does not test semantic recovery; it tests whether the reservoir
can pass consonant-class information that is present in the input features.

\textbf{Mandatory pre-experiment checks} (per Phase~13 protocol):
label function unit tests passed (3 test cases); corpus unique root count
reported prior to any augmentation; permutation test pre-committed to
$n = 1{,}000$; success threshold pre-registered at ARI $> 0.05$,
$p < 0.05$.

\textbf{Results.}

\begin{table}[h]
\centering
\begin{tabular}{lcc}
\toprule
Condition & ARI & $p$ (permutation, $n=1000$) \\
\midrule
Acoustic-only (locus labels) & 0.0494 & $< 0.001$ \\
MLP on 23D acoustic (5-fold CV) & 0.312 $\pm$ 0.021 & $< 0.001$ \\
\bottomrule
\end{tabular}
\caption{Experiment~56 results. The reservoir achieves ARI $= 0.0494$
  ($p < 0.001$) on the locus taxonomy, confirming statistically significant
  acoustic signal processing. The MLP ceiling of 0.312 establishes that
  locus information is present in the features; the reservoir recovers
  approximately 16\% of the available signal.}
\label{tab:exp56}
\end{table}

\textbf{Interpretation.} The acoustic-only ARI of 0.0494 is marginally
below the pre-registered threshold of 0.05 but is highly statistically
significant ($p < 0.001$). Under the Phase~13 decision tree, this
corresponds to Branch~2 (weak but significant signal): the reservoir is
architecturally functional and not the source of the Phase~12 failure.
The MLP ceiling of $0.312$ confirms that substantial locus information
exists in the 23D acoustic features; the reservoir recovers approximately
16\% of this available signal, indicating a representational bottleneck
in the spiking dynamics but not an architectural failure.

The reservoir is not broken. It processes acoustic information above
chance. This rules out the possibility that the Phase~12 negative result
reflects a defective architecture.

\subsection{Experiment 57: Semantic-Locus Correlation}

\textbf{Design.} If the reservoir can recover locus information (Exp~56)
but not semantic information (Phase~12), the failure must be located in
the relationship between the semantic taxonomy and locus structure. We
tested this directly by measuring the Adjusted Rand Index between the
Monier-Williams semantic axis labels and the Pāṇinian locus labels across
the full matched corpus ($N = 1{,}722$ roots).

\textbf{Results.} ARI between semantic taxonomy and locus taxonomy:
$0.0095$ ($p = 0.001$, permutation test $n = 1{,}000$).

\textbf{Interpretation.} The two taxonomies are nearly orthogonal.
A semantic axis label is almost entirely uninformative about locus class,
and vice versa. This measurement is the precise mechanistic diagnosis
of the Phase~12 negative result: the semantic axes (MOT, SEP, CNT, EXP,
TRN) do not align with the articulatory locus classes that acoustic
formants encode. The reservoir cannot recover what the input does not
contain.

\subsection{Mechanistic Characterization}

Together, Experiments~56 and~57 localize the failure precisely:

\begin{enumerate}[noitemsep]
  \item The reservoir recovers locus information from acoustic features
    (Exp~56: ARI $= 0.0494$, $p < 0.001$).
  \item The semantic taxonomy does not correlate with locus structure
    (Exp~57: ARI $= 0.0095$, $p = 0.001$).
  \item Therefore, the semantic taxonomy is not recoverable from
    acoustic features through the locus channel --- the only channel
    the features encode at first order.
\end{enumerate}

This is a stronger result than the Phase~12 negative finding alone. It
establishes that the failure is \emph{task-specific}: the acoustic feature
set is insufficient for the Monier-Williams taxonomy not because the
reservoir fails to process acoustic information, but because the semantic
axes and the acoustic structure are nearly independent at first order.

%% ═══════════════════════════════════════════════════════════════════
\section{Discussion}

\subsection{Why First-Order Formants Are Insufficient}

The Phase~13 diagnosis (Exp~57: semantic-locus ARI $= 0.0095$) provides
a quantitative account of why the Phase~12 negative result was obtained.
The five semantic axes --- Motion, Separation, Containment, Experiential,
and Transfer --- derived from Monier-Williams glosses do not align with
the five Pāṇinian place-of-articulation classes that acoustic formants
encode. Motion roots begin with velars (gam, go), palatals (car, move),
and dentals (dhāv, run) in roughly equal proportion; Containment roots
include labials and velars; Separation roots span dental, retroflex, and
fricative onsets. The semantic-locus ARI of $0.0095$ quantifies this
near-orthogonality directly.

No acoustic feature combination based on first-order formants can
reliably separate these overlapping distributions, because the semantic
axes were not constructed with reference to articulatory structure. The
phonosemantic hypothesis, as operationalized in the DDIN program, required
a consonant-to-axis correspondence that does not exist in this corpus at
this granularity. This is now an empirically measured fact rather than
an inference.

\subsection{Relationship to Earlier Results}

Papers 1 and 2 of the DDIN program reported above-chance ARI on a
150-root benchmark. Those experiments used consonant-class-derived labels.
Under that regime, both baseline and experimental conditions measured
consonant class separation rather than semantic recovery. The +61\%
improvement reported in Paper 2 (static ARI $= 0.037$ vs.\ sequential
ARI $= 0.059$) used the same circular labeling for both conditions; the
relative improvement may reflect a genuine advantage of sequential
encoding for consonant class discrimination, but cannot be interpreted
as evidence for phonosemantic recovery.

Paper 4's ARI $= 0.9758$ was a near-tautological result: labels were
initial consonants, acoustic features were initial consonant formants,
and the reservoir at $N = 2{,}000$ with a 1024D bottleneck produced
near-perfect consonant-class separation. That record has been retracted
from Zenodo.

\subsection{What Remains Valid}

Four contributions from the DDIN program are independent of the labeling
regime and remain valid.

\textbf{The ESL characterization.} The formal proof and empirical
demonstration that LIF-based SNNs with continuous Poisson input collapse
into $\sim$500\,Hz synchronous firing when $\tau_m$ exceeds the critical
threshold (Theorem 1 of \citet{kumar2026esl}) is a mathematical result
about LIF dynamics, not about phonosemantics.

\textbf{BCM homeostasis as ESL resolution.} The AdEx reservoir with
per-neuron BCM thresholds maintains stable asynchronous firing under
language-scale input (2.8 spikes/root, no epileptiform activity across
2,000-root runs).

\textbf{Locus recovery (Phase 13).} The reservoir recovers articulatory
locus information from acoustic features at ARI $= 0.0494$
($p < 0.001$, Exp~56). This partially rehabilitates the sequential ODE
improvement in Paper~2: the +61\% improvement over static embeddings
reflects a genuine sequential encoding advantage for consonant-class
discrimination, even if it cannot be interpreted as semantic recovery.

\textbf{The phonosemantic manifold $\mathcal{M}$.} The theoretical
framework defining phoneme descriptors as four-dimensional articulatory
coordinates and the harmonic coherence metric $H$
\citep{kumar2026phonosemantic} are not invalidated by the failure of
a specific empirical operationalization.

\subsection{What Would Be Required for a Positive Result}

The negative result constrains but does not close the phonosemantic
hypothesis. Three modifications could make the hypothesis testable at
above-chance levels:

\textbf{A consonant-aligned taxonomy.} A semantic taxonomy in which axis
membership correlates with consonant place-of-articulation would make
the acoustic channel genuinely predictive. The Vedic semantic tradition
offers alternative root classifications (e.g., by \textit{kāraka} or
thematic role) that may align more closely with articulatory structure
than the Monier-Williams gloss-derived axes used here.

\textbf{Richer acoustic representations.} Locus equations and
steady-state formants capture place-of-articulation differences between
consonant classes. They do not capture voice onset time, spectral
envelope, prosodic weight, or the acoustic correlates of aspiration and
retroflexion beyond F2. A richer feature set might recover the finer
distinctions required for the semantic taxonomy used here.

\textbf{A different language substrate.} The phonosemantic hypothesis
predicts stronger signal in Sanskrit than in English. A cross-linguistic
experiment measuring acoustic-only ARI across Sanskrit, Hindi, and
English under the same taxonomy would directly test this prediction.
If ARI$_\text{Sanskrit} >$ ARI$_\text{Hindi} >$ ARI$_\text{English}$
even at chance-adjacent levels, the ordering itself is informative.

%% ═══════════════════════════════════════════════════════════════════
\section{Conclusion}

We have presented a systematic negative result with mechanistic
characterization. First-order acoustic formant features, processed through
an AdEx spiking reservoir with BCM homeostasis, do not recover the
Monier-Williams semantic axis structure in the Sanskrit
\textit{Dhātupāṭha} above chance level (ARI $= 0.013$, Phase~12).
Phase~13 locates the failure precisely: the reservoir is architecturally
functional (Exp~56: locus taxonomy ARI $= 0.0494$, $p < 0.001$), but
the target semantic taxonomy is nearly orthogonal to the articulatory
locus structure that acoustic features encode (Exp~57: semantic-locus
ARI $= 0.0095$). The reservoir cannot recover what the input does not
contain.

Earlier positive results in the DDIN program were identified as artifacts
of a label assignment function that derived semantic axes from the same
initial consonant used for acoustic encoding. After correction, no
phonosemantic signal is detectable at the level of first-order formants
for the five-axis Monier-Williams taxonomy. The Speed-1 result
(ARI $= 0.204$ at first phoneme) does not replicate under corrected
labels (ARI $= 0.016$).

Three contributions from the program are unaffected: the formal
characterization of the Epileptiform Synchrony Limit
\citep{kumar2026esl}; the BCM homeostasis architecture that resolves it;
and the Phase~13 finding that the reservoir successfully recovers
articulatory locus information from acoustic features (ARI $= 0.0494$,
$p < 0.001$), which partially rehabilitates the sequential ODE
improvement reported in Paper~2 as a genuine result for consonant-class
discrimination.

This is a complete, mechanistically characterized finding. It specifies:
(1) where the acoustic signal exists --- locus structure, recoverable at
ARI $\approx 0.05$ in the reservoir and ARI $\approx 0.31$ in a
feedforward classifier; (2) where it does not exist --- the
Monier-Williams semantic taxonomy, which is nearly orthogonal to locus;
and (3) what would be required for a positive phonosemantic result --- a
semantic taxonomy whose axes align with articulatory locus structure. We
offer this as a methodological contribution to the literature on
phonosemantics, reservoir computing, and neuromorphic language
processing.

%% ═══════════════════════════════════════════════════════════════════
\section*{Acknowledgements}

This research was conducted independently. The corpus uses the
Monier-Williams Sanskrit-English Dictionary (1899), digitized by the
University of Cologne Sanskrit Lexicon project. No funding was received.
The author thanks the critical engagement from prior review rounds
that precipitated the labeling audit reported here.

\section*{Competing Interests}
The author declares no competing interests.

\section*{Data and Code Availability}
Experiment code, the corrected labeling function, and the full three-condition
ablation results are available at
\href{https://github.com/HmbleCreator}{github.com/HmbleCreator}.

%% ═══════════════════════════════════════════════════════════════════
\begin{thebibliography}{99}

\bibitem{kumar2026phonosemantic}
A.~Kumar.
\newblock Phonosemantic grounding: {S}anskrit as a formalized case of
  motivated sign structure for interpretable {AI}.
\newblock \textit{Zenodo Preprint}, 2026.
\newblock \doi{10.5281/zenodo.19564026}

\bibitem{kumar2026ddin}
A.~Kumar.
\newblock Sequential phonosemantic encoding in {ODE} reservoirs: Breaking
  static baselines and measuring capacity ceilings.
\newblock \textit{Zenodo Preprint}, 2026.
\newblock \doi{10.5281/zenodo.19615067}

\bibitem{kumar2026esl}
A.~Kumar.
\newblock The epileptiform synchrony limit: {E/I} equilibrium as a physical
  prerequisite for intelligence in spiking neural networks.
\newblock \textit{Zenodo Preprint}, 2026.
\newblock \doi{10.5281/zenodo.19602055}

\bibitem{hinton1994}
L.~Hinton, J.~Nichols, and J.~Ohala, editors.
\newblock \textit{Sound Symbolism}.
\newblock Cambridge University Press, 1994.

\bibitem{magnus2001}
M.~Magnus.
\newblock What's in a word? {S}tudies in phonosemantics.
\newblock PhD thesis, Norwegian University of Science and Technology, 2001.

\bibitem{blasi2016}
D.~E. Blasi, S.~Wichmann, H.~Hammarstr{\"o}m, P.~F. Stadler, and
  M.~Christiansen.
\newblock Sound-meaning association biases evidenced across thousands of
  languages.
\newblock \textit{Proceedings of the National Academy of Sciences},
  113(39):10818--10823, 2016.

\bibitem{saussure1916}
F.~de~Saussure.
\newblock \textit{Cours de linguistique g{\'e}n{\'e}rale}.
\newblock Payot, 1916.

\bibitem{cardona1997}
G.~Cardona.
\newblock \textit{P{\={a}}{\d{n}}ini: His Work and Its Traditions}.
\newblock Motilal Banarsidass, 1997.

\bibitem{monierWilliams1899}
M.~Monier-Williams.
\newblock \textit{A Sanskrit-English Dictionary}.
\newblock Oxford University Press, 1899.

\bibitem{sussman1991}
H.~M. Sussman, H.~A. McCaffrey, and S.~A. Matthews.
\newblock An investigation of locus equations as a source of relational
  invariance for stop place categorization.
\newblock \textit{Journal of the Acoustical Society of America},
  90(3):1309--1325, 1991.

\bibitem{goyal2015}
P.~Goyal, A.~Kumar, and A.~Kulkarni.
\newblock A {P}{\={a}}{\d{n}}inian approach to natural language processing.
\newblock In \textit{Proceedings of the Workshop on Computational Linguistics
  for South Asian Languages}, 2015.

\bibitem{brette2005}
R.~Brette and W.~Gerstner.
\newblock Adaptive exponential integrate-and-fire model as an effective
  description of neuronal activity.
\newblock \textit{Journal of Neurophysiology}, 94(5):3637--3642, 2005.

\bibitem{bienenstock1982}
E.~L. Bienenstock, L.~N. Cooper, and P.~W. Munro.
\newblock Theory for the development of neuron selectivity: orientation
  specificity and binocular interaction in visual cortex.
\newblock \textit{Journal of Neuroscience}, 2(1):32--48, 1982.

\bibitem{jaeger2001}
H.~Jaeger.
\newblock The ``echo state'' approach to analysing and training recurrent
  neural networks.
\newblock \textit{GMD Technical Report}, 148, 2001.

\bibitem{maass2002}
W.~Maass, T.~Natschl{\"a}ger, and H.~Markram.
\newblock Real-time computing without stable states: A new framework for
  neural computation based on perturbations.
\newblock \textit{Neural Computation}, 14(11):2531--2560, 2002.

\bibitem{bellman1961}
R.~Bellman.
\newblock \textit{Adaptive Control Processes: A Guided Tour}.
\newblock Princeton University Press, 1961.

\bibitem{dingemanse2015}
M.~Dingemanse, F.~Blasi, G.~Lupyan, M.~H. Christiansen, and P.~Monaghan.
\newblock Arbitrariness, iconicity, and systematicity in language.
\newblock \textit{Trends in Cognitive Sciences}, 19(10):603--615, 2015.

\end{thebibliography}

\end{document}
